\documentclass[11pt,leqno]{amsart}
\usepackage[T1]{fontenc}
\usepackage{lmodern}
\usepackage{amssymb,graphicx,array,booktabs,longtable,listings,xcolor}
\usepackage[hidelinks]{hyperref}
\hypersetup{
  pdftitle={A Log-Sine Analogue of Euler's Constant: The Decay Rate of a Sine Product at Primes},
  pdfauthor={V\'ictor Casa\~na-Giner},
  pdfsubject={A log-sine product at primes, Mellin transforms, and numerical certificates}}
\setlength{\textwidth}{32pc}
\setlength{\textheight}{54pc}
\calclayout
\raggedbottom
\numberwithin{equation}{section}
\newtheorem{theorem}{Theorem}[section]
\newtheorem{proposition}[theorem]{Proposition}
\newtheorem{lemma}[theorem]{Lemma}
\newcommand{\Cg}{C_{\mathrm g}}
\newcommand{\R}{\operatorname{Re}}
\newcommand{\I}{\operatorname{Im}}
\newcommand{\E}{\mathbb{E}}
\newcommand{\decfirst}{7437948618\,3940054420\,3641994634\,5806610252\,6433197291}
\newcommand{\decsecond}{7222960354\,3808135540\,0286772998\,3871757037\,1013522332}
\lstset{language=Python,basicstyle=\ttfamily\scriptsize,columns=fullflexible,
  keepspaces=true,breaklines=true,showstringspaces=false,numberstyle=\tiny,
  numbers=left,numbersep=6pt,xleftmargin=2em,frame=leftline,
  rulecolor=\color{black!22},aboveskip=8pt,belowskip=8pt}
\newcommand{\sourcefigure}[3]{%
  \par\medskip\noindent
  \begin{minipage}{\linewidth}
  \centering
  \includegraphics[width=#1,height=.36\textheight,keepaspectratio]{castor_v5/#2}
  \refstepcounter{figure}\par\smallskip
  {\small\raggedright\textsc{Figure \thefigure.} #3\par}
  \end{minipage}\par\medskip}
\title[A log-sine analogue of Euler's constant]{A Log-Sine Analogue of Euler's Constant: The Decay Rate of a Sine Product at Primes}
\author{V\'ictor Casa\~na-Giner}
\date{November 11, 2026}
\subjclass[2020]{Primary 11M06, 11K65; Secondary 11N37, 44A15}
\keywords{Euler's constant, sine product, fractional parts, Mellin transform, Riemann zeta function, certified computation}
\begin{document}
\begin{abstract}
For prime $p$ we study $P(p)=\prod_{k=2}^{p-1}\sin(\pi p/k)$.
Its logarithmic size is $\log|P(p)|=-\Cg p+O(p^{1/3}\log^2 p)$, where
$\Cg=0.74379486183940054420\ldots$. The rate per factor in the
limit is $e^{-\Cg}=0.47530676524152038801\ldots$. With the
denominators restricted to $k\le\sqrt p$, the logarithmic size is
$-(\log2)\sqrt p+O(p^{1/3}\log^2p)$. The means $1-\gamma$ and $\Cg$ are obtained from two observables of the
same fractional-part distribution. We relate $\Cg$ to the Mellin
transforms $G$ and $A$, derive a geometrically convergent series in
even zeta values, and record an enclosure fixing the first $100$
decimal digits of $\Cg$.
\end{abstract}
\maketitle

\section{Introduction}
For integers $n\ge2$ put
\begin{equation}\label{eq:product}
 P(n)=\prod_{k=2}^{n-1}\sin\left(\frac{\pi n}{k}\right),\qquad
 \widetilde P(n)=\prod_{2\le k\le\sqrt n}\sin\left(\frac{\pi n}{k}\right).
\end{equation}
The empty products are $P(2)$, $\widetilde P(2)$ and $\widetilde P(3)$,
each with value $1$ by convention. A factor vanishes exactly when $k$ divides
$n$, and every composite $n$ has a divisor in $[2,\sqrt n]$.
Thus each product is nonzero exactly when $n$ is prime. For primes $p$ the logarithm of the absolute
value of the product has the asymptotic form
\begin{equation}\label{eq:decay}
 \log|P(p)|=-\Cg p+O(p^{1/3}\log^2p),
 \qquad \Cg=0.74379486183940054420\ldots .
\end{equation}
Consequently $|P(p)|^{1/p}\to e^{-\Cg}=0.47530676524152038801\ldots$.
The logarithmic decay rate is $\Cg$, and $e^{-\Cg}$ is the
exponential base. For the product truncated at the square root,
\begin{equation}\label{eq:shortdecay}
 \log|\widetilde P(p)|=-(\log2)\sqrt p+O(p^{1/3}\log^2p),
 \qquad |\widetilde P(p)|^{1/\sqrt p}\longrightarrow\frac12 .
\end{equation}
The constants in \eqref{eq:decay} and \eqref{eq:shortdecay} arise
from different distributions of the fractional parts.

Let $U$ be uniform on $(0,1)$ and put $X=\{1/U\}$. This is the
limiting distribution of the fractional parts $\{n/k\}$ for the
full product $P(n)$. The mean of $X$ is $1-\gamma$, where $\gamma$ is
Euler's constant, and the mean of $-\log(\sin\pi X)$ is $\Cg$. These
two means occur in a single complex expectation:
\begin{equation}\label{eq:eulerlink}
 \E\operatorname{Log}(1-e^{2\pi iX})
   =(\log 2-\Cg)+i\pi\left(\frac12-\gamma\right).
\end{equation}
Here $\operatorname{Log}$ is the principal logarithm. This identity is
the sense in which $\Cg$ is a log-sine analogue of Euler's constant.

For real $u\notin\mathbb Z$ put
\begin{equation}\label{eq:h}
h(u)=\log|2\sin(\pi u)| .
\end{equation}
The transform of $h$ on $[1,\infty)$ is
\begin{equation}\label{eq:G}
G(s)=\int_1^\infty h(u)\,u^{-s-1}\,du\qquad(\operatorname{Re}s>0).
\end{equation}
There is also a parallel at the level of Mellin transforms. Put
\[
 \mathcal M_{\mathrm{frac}}(s)=\int_1^\infty\{u\}\,u^{-s-1}\,du
 \qquad(\R s>0).
\]
For $\R s>1$, integrating $\lfloor u\rfloor=\sum_{n\le u}1$
termwise gives
\[
 \zeta(s)=s\int_1^\infty\lfloor u\rfloor u^{-s-1}\,du
         =\frac{s}{s-1}-s\mathcal M_{\mathrm{frac}}(s).
\]
The integral defining $\mathcal M_{\mathrm{frac}}$ is holomorphic on
$\R s>0$ because $\{u\}$ is bounded and the integral converges locally
uniformly there. The formula thus continues $\zeta$ meromorphically
to that half-plane. Comparing constant terms at $s=1$
with the Laurent expansion of $\zeta$ \cite[(25.2.4)--(25.2.5)]{dlmf}
gives $\mathcal M_{\mathrm{frac}}(1)=1-\gamma$. Likewise,
$\log(2)/s-G(s)$ is the transform of $-\log|\sin(\pi u)|$
on $[1,\infty)$ and takes the value $\Cg$ at $s=1$ by
\eqref{eq:Cint}. Both values are obtained with the weight $u^{-2}$.

For $0<\operatorname{Re}s<2$ let
\begin{equation}\label{eq:A}
A(s)=\frac{\log(2\pi)}s+\frac1{s^2}-\int_0^1\log\frac{\sin(\pi u)}{\pi u}\,u^{-s-1}\,du .
\end{equation}
In particular $\Cg=\log2-G(1)$. As an identity between meromorphic
functions,
\begin{equation}\label{eq:GAzeta}
G(s)-A(s)=\frac\pi s\cot\Bigl(\frac{\pi s}2\Bigr)\zeta(s) .
\end{equation}
For non-real $s$, equality $G(s)=A(s)$ is equivalent to $\zeta(s)=0$,
with the same multiplicities. Thus if every non-real zero of
$\zeta$ has real part $\tfrac12$, the non-real solutions of $G=A$ are
of the form $\tfrac12+it$ with $t\ne0$. The real solutions are the odd
integers other than $1$. At $s=1$ the zero of the cotangent cancels the
pole of $\zeta$, giving $G(1)-A(1)=-\pi^2/2$.

\section{Notation and constants}
Throughout, $s=\sigma+it$ with $\sigma,t$ real. The index $n$ is an
integer, $p$ is a prime, and $k$ and $m$ are summation indices.
The fractional part of $x$ is $\{x\}=x-\lfloor x\rfloor$.
We use $y$ for a variable in $(0,1)$ in the fractional-part
distribution. All real logarithms are natural logarithms.
\begin{center}
 \small
 \begin{tabular}{@{}lll@{}}
 \toprule
 Symbol & Meaning & Initial decimals\\
 \midrule
 $\gamma$ & Euler's constant & $0.57721566490153286060\ldots$\\
 $\Cg$ & Logarithmic decay rate of $|P(p)|$ & $0.74379486183940054420\ldots$\\
 $e^{-\Cg}$ & Exponential base of $|P(p)|$ & $0.47530676524152038801\ldots$\\
 $\log2$ & Logarithmic decay rate of $|\widetilde P(p)|$ & $0.69314718055994530941\ldots$\\
 $G(1)$ & $\log 2-\Cg$ & $-0.05064768127945523478\ldots$\\
 $A(1)$ & $G(1)+\pi^2/2$ & $4.88415451926522407463\ldots$\\
 \bottomrule
 \end{tabular}
 \par\smallskip
 \refstepcounter{table}
 \textsc{Table \thetable.} The constants used in the product and its Mellin representation.
\end{center}


{\small
\setlength{\LTleft}{0pt}
\setlength{\LTright}{0pt}
\begin{longtable}{@{}>{\raggedright\arraybackslash}p{.37\linewidth}
  >{\raggedright\arraybackslash}p{\dimexpr.63\linewidth-2\tabcolsep\relax}@{}}
\caption{Notation used in the proofs and appendices.}\label{tab:symbols}\\
\toprule
Symbol & Meaning\\
\midrule
\endfirsthead
\toprule
Symbol & Meaning\\
\midrule
\endhead
\bottomrule
\endfoot
$n,p,k,m,j,b,l,e_q$ & Integer arguments and indices. In the divisor bound $q$ is prime and $e_q$ is its exponent in $n$.\\
$P,\widetilde P,d(n),Y$ & Full product, square-root product, divisor count, and cutoff in the counting estimate.\\
$U,X,y,y_k$ & Uniform variable on $(0,1)$, its image $\{1/U\}$, a fractional-part variable, and $\{n/k\}$.\\
$w,F$ & Density of $X$ and its distribution function.\\
$\mu,\nu,W,T,D_\mu,\tau$ & Sample measure, signed discrepancy measure, comparison density, comparison mass, uniform discrepancy, and counting threshold.\\
$f,f_\delta,\delta,I,m_0$ & Logarithmic observable, its clipped form, endpoint cutoff, uniform integral of $f$, and a finite sum or product endpoint.\\
$\gamma,\Cg,e^{-\Cg},\log2$ & Euler's constant, the full-product decay rate, its exponential base, and the square-root decay rate.\\
$s,\sigma,t,i,u,v,x,z,\xi$ & Complex argument, its real and imaginary coordinates, the imaginary unit, and local integration and function arguments. In the principal-value calculation $z=s/2$.\\
$\log,\operatorname{Log},\operatorname{PV}$ & Real natural logarithm, principal complex logarithm, and Cauchy principal value.\\
$h,\ell,G,A$ & Log-sine kernel, regularized logarithm, and the two functions in \eqref{eq:G} and \eqref{eq:A}.\\
$\mathcal M_{\mathrm{frac}}$ & Mellin transform of the fractional part on $[1,\infty)$.\\
$\zeta,\Gamma,\Sigma_{2k},J(s),\psi$ & Zeta and gamma functions, even zeta values, the scalar Mellin integral, and the digamma function.\\
$N_C,\Cg^{[N_C]},B_{2k}$ & Number of terms in the series for $\Cg$, its partial value, and the even Bernoulli numbers.\\
$M_C,a,R_{N_C,M_C}$ & Cutoff in the elementary series, the first omitted index $a=M_C+1$, and the remainder of the accelerated formula.\\
$D_{100},L_{180},U_{180}$ & Integer formed by the first 100 fractional digits, and the endpoints enclosing the 180-term partial value.\\
$\mathcal K,\mathcal K_\alpha,H$ & Periodic kernel, its Abel regularization with $0<\alpha<1$, and its Mellin integral.\\
$\mathsf M,\mathsf N,\mathsf L,x_b$ & Numbers of harmonics, periods, and midpoint nodes in the Rouch\'e certificate, and its midpoint nodes.\\
$\mathcal K_{\mathsf M},H_{\mathsf M,\mathsf N},
 H_{\mathsf M,\mathsf N,\mathsf L}$ & Kernel cutoff, cutoff integral, and midpoint approximation.\\
$B_0,B_1,B_2,Z,E_0,E_1$ & Bounds for the kernel and its derivatives, an upper bound for a sum of inverse powers, and the errors for $H$ and $H'$.\\
$s_0,\varsigma,\eta_0,\eta_1$ & Disc center and radius, and the approximations to $H(s_0)$ and $H'(s_0)$.\\
$Q_0,\mathcal Q,\mathfrak r_k,c_k,g_\psi$ & Digamma remainder series, upper bound for $\R(sA(s))$, $\R(s/(2k-s))$, $\Sigma_{2k}-1$, and the nonnegative digamma comparison integrand.\\
$\varepsilon,\rho$ & Exponent in the divisor bound and cutoff for the principal value.\\
$u_0,\theta,\lambda,\chi$ & Tail endpoint, inverse-power exponent, local integration variable, and midpoint-rule integrand.\\
$g_{jk},r,\beta$ & Greatest common divisor $\gcd(j,k)$, positive integer multiple, and frequency in the Gram calculation.\\
$z_b^G,z_b^A,n_G,n_A,\mathcal E_G,\mathcal E_A$ & Lists of computed zeros, their lengths, and the exponential sums.\\
\end{longtable}
}

\section{The distribution and the arithmetic product}
Splitting $1/U$ into intervals $[m,m+1)$ gives the density
\begin{equation}\label{eq:density}
 w(y)=\sum_{m=1}^{\infty}\frac{1}{(m+y)^2},\qquad 0<y<1 .
\end{equation}
Indeed the change of variables $u=1/(m+y)$ on each interval of
$U$ contributes $(m+y)^{-2}\,dy$. The sum of the integrals of
these terms is $1$. This distribution for $\{n/k\}$ is treated
by Saffari and Vaughan \cite[Corollary 1.1]{saffariVaughan}
and Bensimhoun \cite[Examples 3--4]{bensimhoun}. Termwise integration of
$yw(y)$ up to the summand $m=m_0$ gives
\[
 \sum_{m=1}^{m_0}\left(\log\frac{m+1}{m}-\frac1{m+1}\right)
 =1+\log(m_0+1)-\sum_{m=1}^{m_0+1}\frac1m .
\]
Its limit is $1-\gamma$. Integrating the log-sine observable gives
\begin{equation}\label{eq:Cint}
 \Cg=-\int_0^1\log(\sin\pi y)\,w(y)\,dy
     =-\int_1^\infty\frac{\log|\sin(\pi u)|}{u^2}\,du .
\end{equation}
The first integral converges because $w$ is bounded and the endpoint
singularities are logarithmic. The second equality follows by splitting
the integral on $[1,\infty)$ into unit intervals. For $0<y<1$ the
identity $1-e^{2\pi iy}=2\sin(\pi y)e^{i\pi(y-1/2)}$ has argument
in $(-\pi/2,\pi/2)$. Taking the mean of its principal logarithm gives
\eqref{eq:eulerlink}.

\begin{lemma}[Discrepancy and the logarithmic endpoints]\label{lem:arithmetic}
Let $n\ge3$, and let $\mu$ count a finite list of points in
$[1/n,1-1/n]$, with multiplicities. Let $W$ be a bounded nonnegative
density of integral $1$ on $[0,1]$, and let $T>0$. Put
\[
 D_\mu=\sup_{0\le\tau\le1}\left|\mu([0,\tau])-T\int_0^\tau W(y)\,dy\right|.
\]
For $f(y)=\log(\sin\pi y)$ and $1/n\le\delta<1/2$,
\[
 \int f\,d\mu=T\int_0^1 f(y)W(y)\,dy+O((D_\mu+T\delta)\log n),
\]
with an implied constant depending only on an upper bound for $W$.
\end{lemma}

\begin{proof}
Define $f_\delta=f$ on $[\delta,1-\delta]$ and extend it constantly
to each endpoint of $[0,1]$. Since
$f_\delta(1)=f(1-\delta)=\log(\sin\pi\delta)$, the inequality
$\sin(\pi y)\ge2\min(y,1-y)$ and the monotonicity of $f$ on each
side of $1/2$ give
\[
 |f_\delta(1)|+\int_0^1|f_\delta'(y)|\,dy=O(\log n).
\]
For $\nu=\mu-TW(y)\,dy$, insert
$f_\delta(y)=f_\delta(1)-\int_y^1 f_\delta'(\tau)\,d\tau$ in
$\int f_\delta\,d\nu$ and interchange the finite-measure integrals.
The result is bounded in modulus by
$D_\mu\bigl(|f_\delta(1)|+\int_0^1|f_\delta'|\bigr)$.
The two endpoint intervals contain $O(T\delta+D_\mu)$ sample points
by the definition of $D_\mu$ and boundedness of $W$. Both $f$ and
$f_\delta$ have size $O(\log n)$ at those points, since the
sample lies in $[1/n,1-1/n]$ and $\delta\ge1/n$. On the integral
side,
$\int_0^\delta|\log(2y)|\,dy=O(\delta(1+|\log\delta|))$,
with the same bound at the other endpoint by reflection. Thus
replacing $f_\delta$ by $f$ costs $O((T\delta+D_\mu)\log n)$ in the
sample sum and $O(T\delta\log n)$ in the integral against
$TW(y)\,dy$. These are the claimed bounds.
\end{proof}

\begin{theorem}[Arithmetic sampling]\label{thm:arithmetic}
Let $d(n)$ be the number of positive divisors of $n$. For every integer
$n\ge2$, with absolute implied constants,
\begin{align}
 \sum_{\substack{2\le k\le n-1\\k\nmid n}}
    \log\left|\sin\frac{\pi n}{k}\right|
  &=-\Cg n+O\bigl(n^{1/3}\log^2n+d(n)\log n\bigr),
   \label{eq:arithmetic}\\
 \sum_{\substack{2\le k\le\sqrt n\\k\nmid n}}
    \log\left|\sin\frac{\pi n}{k}\right|
  &=-(\log2)\sqrt n+O\bigl(n^{1/3}\log^2n+d(n)\log n\bigr).
   \label{eq:shortarithmetic}
\end{align}
\end{theorem}

\begin{proof}
For the full range, let
$F(\tau)=\int_0^\tau w(y)\,dy=\sum_{m\ge1}\tau/[m(m+\tau)]$
for $0\le\tau\le1$.
Saffari and Vaughan \cite[Corollary 1.1, equation (1.9)]{saffariVaughan}
give
\[
 \#\{1\le k\le n:\{n/k\}<\tau\}=nF(\tau)+O(n^{1/3}\log n)
 \qquad(0\le\tau\le1)
\]
uniformly in $\tau$. The points at zero are the $d(n)$ divisors of $n$,
including $1$ and $n$. Passing to right limits in the counting
estimate and removing the divisors gives
$D_\mu=O(n^{1/3}\log n+d(n))$ for $T=n$ and $W=w$ in
Lemma~\ref{lem:arithmetic}. The endpoint $\tau=1$ follows from the
count of the remaining $n-d(n)$ points.

For the square-root range, Saffari and Vaughan
\cite[Corollary 1.3, equation (1.11)]{saffariVaughan} give, for
$Y=o(n)$, the uniform bound
\[
 \#\{1\le k\le Y:\{n/k\}<\tau\}
 =Y\tau+O(Y^2/n+n^{1/3}\log n)\qquad(0\le\tau\le1).
\]
With $Y=\sqrt n$ the first error term is $1$. As before, taking
right limits and removing the divisors, including $1$, gives
$D_\mu=O(n^{1/3}\log n+d(n))$, this time for $T=\sqrt n$ and $W=1$.
At $\tau=1$ the remaining count is
\[
 \mu([0,1])=\lfloor\sqrt n\rfloor-\#\{k:1\le k\le\sqrt n,\ k\mid n\}.
\]
Thus $|\mu([0,1])-\sqrt n|\le d(n)+1$.
The bounds also give the distributional
limits in each range, since $D_\mu/T\to0$. For the latter assertion
one may use the elementary divisor bound $d(n)=O_\varepsilon(n^\varepsilon)$
for any $\varepsilon>0$. Indeed, if $n=\prod_q q^{e_q}$, then
$e_q+1\le 2^{e_q}\le q^{\varepsilon e_q}$ for $q\ge2^{1/\varepsilon}$.
For each of the finitely many smaller primes $q$, the ratio
$(e_q+1)/q^{\varepsilon e_q}$ is bounded over integers $e_q\ge0$.
Multiplication gives the bound.

In either range $y_k=\{n/k\}$ for $k\nmid n$ satisfies
$\min(y_k,1-y_k)\ge1/k\ge1/n$. Writing
$n/k=\lfloor n/k\rfloor+y_k$ gives
$\sin(\pi n/k)=(-1)^{\lfloor n/k\rfloor}\sin(\pi y_k)$, and hence
$\log|\sin(\pi n/k)|=\log(\sin(\pi y_k))$ because $0<y_k<1$.
Use Lemma~\ref{lem:arithmetic}
with $\delta=n^{-2/3}$ for $T=n$, and with $\delta=n^{-1/6}$ for
$T=\sqrt n$. In both cases $T\delta=n^{1/3}$. The finitely
many $n$ for which the prescribed $\delta$ is at least $1/2$
are covered by enlarging the implied constant. The full-range
integral is $-\Cg$ by \eqref{eq:Cint}. In the square-root range
write $I=\int_0^1\log(\sin\pi y)\,dy$. By a shift of half a
period, $\int_0^1\log|\cos(\pi y)|\,dy=I$. The double-angle
formula then gives $2I=I-\log2$, because
$\int_0^1\log|\sin(2\pi y)|\,dy=I$. This gives $I=-\log2$
and proves \eqref{eq:shortarithmetic}.
\end{proof}

For $n=p$ prime, $d(p)=2$ and neither product loses any factors
in the respective sums. Since $2\log p=O(p^{1/3}\log^2p)$,
equation \eqref{eq:arithmetic} gives \eqref{eq:decay}, while
\eqref{eq:shortarithmetic} gives the logarithmic estimate in
\eqref{eq:shortdecay}. Division by
$\sqrt p$ and exponentiation give the limit $1/2$ in
\eqref{eq:shortdecay}.

\section{The Mellin representation and the zeta series}
For real $u>0$ with $u\notin\mathbb Z$ put
\begin{equation}\label{eq:ell}
\ell(u)=\log\frac{\pi u}{|\sin(\pi u)|} .
\end{equation}
By \eqref{eq:h}, $h(u)-\log(2\pi u)=-\ell(u)$. The definitions
\eqref{eq:G} and \eqref{eq:A} then give
\begin{equation}\label{eq:GAell}
G(s)-A(s)=-\int_0^\infty\ell(u)\,u^{-s-1}\,du\qquad(0<\operatorname{Re}s<2).
\end{equation}
To see this, split the integral at $1$ and use
$\int_1^\infty\log(2\pi u)u^{-s-1}\,du=\log(2\pi)/s+1/s^2$.
Near zero the logarithm is $O(u^2)$. At infinity it is the sum
of a logarithm and a periodic function with integrable singularities.
More explicitly, $\int_m^{m+1}|h(u)|\,du=\int_0^1|h(u)|\,du<\infty$.
These bounds give absolute convergence, locally uniformly for
$0<\R s<2$, and show that \eqref{eq:G} is holomorphic for $\R s>0$.
In the notation of Adam \cite[Theorem 2.2, definition of $I_s$]{adam},
the specialization $(1,0,\tfrac14,1)$ gives, for $\R s>0$,
\[
 I_s(1,0,\tfrac14,1)
   =-\frac12\int_1^\infty\log\bigl(4\sin^2(\pi u)\bigr)u^{-s-1}\,du
   =-G(s).
\]
In particular, $\Cg=\log2+I_1(1,0,\tfrac14,1)$.

\begin{theorem}[Continuation and Mellin identity]\label{thm:mellin}
The function $A$ extends meromorphically to $\mathbb C$ by the series
\begin{equation}\label{eq:Aseries}
A(s)=\frac{\log(2\pi)}{s}+\frac1{s^2}+\sum_{k\ge1}\frac{\Sigma_{2k}}{k\,(2k-s)} ,
\end{equation}
where $\Sigma_{2k}=\sum_{m\ge1}m^{-2k}=\zeta(2k)$. The function $G$
extends to an entire function and \eqref{eq:GAzeta} holds meromorphically
on $\mathbb C$.
\end{theorem}

\begin{proof}
Euler's product for the sine \cite[(4.22.1)]{dlmf} gives, for $0<u<1$,
\[
 -\log\frac{\sin(\pi u)}{\pi u}
 =\sum_{m\ge1}\sum_{k\ge1}\frac{u^{2k}}{k m^{2k}}
 =\sum_{k\ge1}\frac{\Sigma_{2k}}{k}u^{2k}.
\]
For $\sigma<2$ the sum of the integrals of the absolute values
after multiplication by $u^{-s-1}$ is
$\sum_{k\ge1}\Sigma_{2k}/[k(2k-\sigma)]<\infty$. Integration
therefore proves \eqref{eq:Aseries}. The large-$k$ terms are
$O(k^{-2})$ uniformly on each compact set, which gives the claimed
meromorphic continuation. Its only poles are a double pole at $0$
and simple poles at the positive even integers.

For positive nonintegral $u$, the same sine product gives
$-\ell(u)=\sum_{m\ge1}\log|1-u^2/m^2|$. If $1<\sigma<2$, the
sum of the weighted absolute integrals on $(0,\infty)$ is
\[
 \sum_{m\ge1}m^{-\sigma}
 \int_0^\infty |\log|1-v^2||\,v^{-\sigma-1}\,dv<\infty.
\]
The inner integral is finite by the $O(v^2)$ behavior at zero,
the logarithmic singularity at $1$, and $O(\log v)$ growth at
infinity. Termwise integration in \eqref{eq:GAell}, with $u=mv$,
now gives $G(s)-A(s)=\zeta(s)J(s)$, where
\[
 J(s)=\int_0^\infty\log|1-v^2|\,v^{-s-1}\,dv
      =\frac12\int_0^\infty\log|1-\xi|\,\xi^{-s/2-1}\,d\xi.
\]
Write $z=s/2$. Integration by parts on the intervals separated by
$1-\rho$ and $1+\rho$ gives
$J(s)=s^{-1}\operatorname{PV}\int_0^\infty \xi^{-z}/(\xi-1)\,d\xi$.
The boundary terms at zero and infinity vanish for $0<\R z<1$,
and the paired terms at $1\pm\rho$ are
$O(\rho|\log\rho|)$. Substitution $\xi=1/y$ above $1$
changes the cutoff there to $1/(1+\rho)$. Replacing it by
$1-\rho$ changes the integral by $O(\rho)$, since the
endpoints are $O(\rho^2)$ apart and the integrand there is
$O(\rho^{-1})$. The principal value is
\[
 \int_0^1\frac{y^{z-1}-y^{-z}}{1-y}\,dy
 =\sum_{m=0}^\infty
       \left(\frac1{m+z}-\frac1{m+1-z}\right)
 =\pi\cot(\pi z).
\]
The geometric-series integration is justified by the integrability
of the modulus of the difference divided by $1-y$, since the
numerator vanishes at $1$. The last equality is the partial
fraction expansion of the cotangent \cite[(4.22.3)]{dlmf}.
This proves \eqref{eq:GAzeta} on $1<\R s<2$.

Define the continuation of $G$ by adding the right side of
\eqref{eq:GAzeta} to \eqref{eq:Aseries}. At a positive even
integer $2k$, the residues of the two summands are
$-\zeta(2k)/k$ and $\zeta(2k)/k$. At a negative even integer
the zero of $\zeta$ \cite[(25.6.4)]{dlmf} cancels the cotangent pole. The factor
$\cot(\pi s/2)$ cancels the pole of $\zeta$ at $1$.
At zero,
$(\pi/s)\cot(\pi s/2)=2/s^2+O(1)$ and
$\zeta(0)=-1/2$, $\zeta'(0)=-\log(2\pi)/2$
\cite[(25.6.1) and (25.6.11)]{dlmf}.
The two principal parts at zero cancel. There are no other
possible poles. This entire function agrees with the integral
definition of $G$ on $1<\R s<2$. Since both functions are holomorphic
on the connected half-plane $\R s>0$, the identity theorem shows
that they agree throughout $\R s>0$.
\end{proof}

\begin{proposition}\label{prop:zeros}
For every non-real $s$, $G(s)=A(s)$ if and only if $\zeta(s)=0$, and the multiplicity of $s$ as a zero of $G-A$ equals its multiplicity as a zero of $\zeta$.
The real zeros of $G-A$ are simple and occur at the odd integers
different from $1$.
\end{proposition}
\begin{proof}
For non-real $s$ the factor $(\pi/s)\cot(\pi s/2)$ in
\eqref{eq:GAzeta} is holomorphic and nonzero, so it preserves zeros
and their multiplicities. For real $s<0$, the functional equation
\cite[(25.4.2)]{dlmf} gives, after cancellation of the sine factor,
\[
 G(s)-A(s)=\frac{(2\pi)^s}{s}\cos\frac{\pi s}{2}\,
           \Gamma(1-s)\zeta(1-s).
\]
The equality holds at the negative even integers by continuation.
Since $\Gamma(1-s)$ and $\zeta(1-s)$ are positive and finite on
$s<0$, the only zeros there are the simple cosine zeros at the
negative odd integers.

For $0<s<1$, \eqref{eq:GAell} gives $G(s)-A(s)<0$, since
$|\sin(\pi u)|<\pi u$ for $u>0$. For $s>1$, the Dirichlet series
gives $\zeta(s)>0$. The cotangent factor then gives simple zeros at
$3,5,\ldots$ and poles at the positive even integers. This covers
all the real nonintegers as well. At $0$ there is a double pole,
since $\zeta(0)=-1/2$ and $(\pi/s)\cot(\pi s/2)=2/s^2+O(1)$. At $1$,
$(s-1)\zeta(s)\to1$ and $\cot(\pi s/2)=-\pi(s-1)/2+O((s-1)^3)$.
Hence the value at $1$ is $-\pi^2/2$.
\end{proof}

\begin{theorem}\label{thm:series}
The constant $\Cg$ satisfies
\begin{equation}\label{eq:series}
 \Cg=\frac{\pi^2}{2}-\log(2\pi)-1
       -2\sum_{k=1}^{\infty}\frac{\zeta(2k)-1}{2k-1}.
\end{equation}
If $\Cg^{[N_C]}$ denotes the right side with $1\le k\le N_C$, then
\begin{equation}\label{eq:tail}
 \frac{2\bigl(\zeta(2N_C+2)-1\bigr)}{2N_C+1}
 \le \Cg^{[N_C]}-\Cg
 \le \frac{2\cdot4^{-N_C}}{2N_C+1}\qquad(N_C\ge1).
\end{equation}
\end{theorem}

\begin{proof}
Equation \eqref{eq:Cint} gives $\Cg=\log2-G(1)$.
The value of \eqref{eq:GAzeta} at $1$ gives
$G(1)=A(1)-\pi^2/2$. Substitution of \eqref{eq:Aseries}
therefore yields
\[
 \Cg=\frac{\pi^2}{2}-\log\pi-1
         -\sum_{k\ge1}\frac{\zeta(2k)}{k(2k-1)} .
\]
The finite product
$\prod_{m=2}^{m_0}(1-m^{-2})=(m_0+1)/(2m_0)$ gives
$\sum_{k\ge1}(\zeta(2k)-1)/k=\log2$ by summing the logarithmic
series, whose terms are nonnegative. Also
$\sum_{k\ge1}1/[k(2k-1)]=2\log2$ by the alternating harmonic
series. Using $1/[k(2k-1)]=2/(2k-1)-1/k$ gives \eqref{eq:series}.
For the remainder,
\[
 0<\zeta(2k)-1\le 4^{-k}+\int_2^\infty v^{-2k}\,dv\le 3\cdot 4^{-k}.
\]
For $k>N_C$ we have $2k-1\ge2N_C+1$ and
$\sum_{k>N_C}3\cdot4^{-k}=4^{-N_C}$, which gives the upper bound.
Keeping the first positive term of the remainder gives the lower
bound. Since $\zeta(2N_C+2)-1\ge4^{-(N_C+1)}$, the remainder has
order $4^{-N_C}/N_C$ as $N_C\to\infty$.
\end{proof}

\begin{proposition}[Elementary series and acceleration]\label{prop:hurwitz}
The constant also satisfies
\begin{equation}\label{eq:elementary-series}
 \Cg=\frac{\pi^2}{2}-\log(2\pi)-1
     -\sum_{m\ge2}\frac{1}{m}\log\frac{m+1}{m-1}.
\end{equation}
For any integer $M_C\ge2$, put $a=M_C+1$. Then
\begin{equation}\label{eq:hurwitz-series}
 \begin{aligned}
 \Cg={}&\frac{\pi^2}{2}-\log(2\pi)-1
        -\sum_{m=2}^{M_C}\frac{1}{m}\log\frac{m+1}{m-1}\\
       &{}-2\sum_{k\ge1}\frac{\zeta(2k,a)}{2k-1},
 \end{aligned}
\end{equation}
where $\zeta(s,a)$ is the Hurwitz zeta function. Truncating the last
sum at $k=N_C$ gives an overestimate with remainder $R_{N_C,M_C}$
satisfying
\begin{equation}\label{eq:hurwitz-tail}
 \frac{2\zeta(2N_C+2,a)}{2N_C+1}
 \le R_{N_C,M_C}\le
 \frac{2\zeta(2N_C+2,a)}{(2N_C+1)(1-a^{-2})}.
\end{equation}
\end{proposition}

\begin{proof}
All terms in the double series below are nonnegative. Tonelli's
theorem and the series for $\operatorname{artanh}$ give
\[
 \begin{aligned}
 2\sum_{k\ge1}\frac{\zeta(2k)-1}{2k-1}
 &=2\sum_{m\ge2}\sum_{k\ge1}\frac{m^{-2k}}{2k-1}\\
 &=2\sum_{m\ge2}\frac1m\operatorname{artanh}\frac1m
   =\sum_{m\ge2}\frac1m\log\frac{m+1}{m-1}.
 \end{aligned}
\]
The summand is $2m^{-2}+O(m^{-4})$. Keeping the terms
$2\le m\le M_C$ in this form and using
$\zeta(2k,a)=\sum_{m\ge a}m^{-2k}$ \cite[(25.11.1)]{dlmf}
for the remaining terms proves \eqref{eq:hurwitz-series}.
The first omitted term gives the lower bound in
\eqref{eq:hurwitz-tail}. For the upper bound use
$\zeta(2k+2,a)\le a^{-2}\zeta(2k,a)$ and sum a geometric
series, bounding all the denominators below by $2N_C+1$.
Finally, for $k\ge1$, the integral test gives
\[
 a^{-2k}\le\zeta(2k,a)\le a^{-2k}
               \left(1+\frac{a}{2k-1}\right).
\]
Thus for each fixed $M_C$ the remainder has order
$(M_C+1)^{-2N_C}/N_C$ as $N_C\to\infty$.
\end{proof}

\begin{thebibliography}{99}
\bibitem{adam}
A.~Adam, \emph{Generalized hypergeometric expansion related to the Hurwitz
zeta function}, arXiv:2009.07134v1, 2020.
\href{https://doi.org/10.48550/arXiv.2009.07134}{doi:10.48550/arXiv.2009.07134}.

\bibitem{bensimhoun}
M.~Bensimhoun, \emph{A framework for asymptotic limit problems of
probabilistic nature}, arXiv:2404.03513v2, 2024.
\href{https://doi.org/10.48550/arXiv.2404.03513}{doi:10.48550/arXiv.2404.03513}.

\bibitem{dlmf}
NIST Digital Library of Mathematical Functions,
\href{https://dlmf.nist.gov/4}{Chapter 4}, \href{https://dlmf.nist.gov/5}{Chapter 5}
and \href{https://dlmf.nist.gov/25}{Chapter 25}.

\bibitem{johansson}
F.~Johansson, \emph{Arb: efficient arbitrary-precision midpoint-radius
interval arithmetic}, IEEE Trans. Comput. \textbf{66} (2017), 1281--1292.
\href{https://doi.org/10.1109/TC.2017.2690633}{doi:10.1109/TC.2017.2690633}.

\bibitem{saffariVaughan}
B.~Saffari and R.~C.~Vaughan, \emph{On the fractional parts of $x/n$ and
related sequences. II}, Ann. Inst. Fourier (Grenoble) \textbf{27}
(1977), no.~2, 1--30.
\href{https://doi.org/10.5802/aif.649}{doi:10.5802/aif.649}.
\end{thebibliography}

\clearpage
\appendix
\section{Numerical certificates}\label{app:cert}
\subsection{Decimal enclosure for \texorpdfstring{$\Cg$}{C-g}}
The computation retained $N_C=180$ terms of \eqref{eq:series} in
150-decimal interval arithmetic. The even zeta values were enclosed
using Euler's formula \cite[(25.6.2)]{dlmf},
\[
 \zeta(2k)=\frac{|B_{2k}|(2\pi)^{2k}}{2(2k)!}
\]
with exact rational Bernoulli numbers $B_{2k}$. If $[L_{180},U_{180}]$
encloses the partial sum, the computation used the weaker error
bound $2\cdot4^{-180}$ implied by \eqref{eq:tail}. Thus
\[
 \Cg\in[L_{180}-2\cdot4^{-180},U_{180}].
\]
The two endpoints, multiplied by $10^{100}$, have the same integer
part $D_{100}$, whose digits are printed below. Thus
$D_{100}\,10^{-100}\le\Cg<(D_{100}+1)10^{-100}$.
\begin{equation}\label{eq:100digits}
 \begin{aligned}
  \Cg={}&0.\decfirst\\[-2pt]
         &\phantom{0.}\decsecond\ldots .
 \end{aligned}
\end{equation}
The recorded result and the calculation script are
\texttt{certificate\_Cg\_v1.0.txt} and \texttt{verify\_Cg\_v1.0.py} in the
accompanying files.

\subsection{Kernel and Rouch\'e test for \texorpdfstring{$G$}{G}}\label{sec:rouche}
Put
\begin{equation}\label{eq:K}
 \mathcal K(u)=\sum_{m\ge1}\frac{1-\cos(2\pi m u)}{m^3}.
\end{equation}
Its Mellin integral is
\begin{equation}\label{eq:H}
 H(s)=\int_1^\infty\mathcal K(u)u^{-s-3}\,du\quad(\R s>-2).
\end{equation}
The kernel series converges uniformly, and
$0\le\mathcal K\le2\zeta(3)<5/2$ by the integral test starting
at $m=2$. In particular $H$ is holomorphic on $\R s>-2$.

To relate $H$ to $G$, insert a factor $\alpha^m$ in the kernel series
for $0<\alpha<1$ and denote the resulting kernel by $\mathcal K_\alpha$.
The logarithmic series at $\alpha e^{2\pi iu}$ gives
$\mathcal K_\alpha''(u)=-4\pi^2\log|1-\alpha e^{2\pi iu}|$. Integration by
parts twice against $u^{-s-1}$ on $[1,\infty)$ for $\R s>0$
has zero boundary terms, since $\mathcal K_\alpha(1)=\mathcal K_\alpha'(1)=0$
and the kernel and its derivative are bounded. To let $\alpha\uparrow1$,
use uniform convergence of $\mathcal K_\alpha$ to $\mathcal K$ and,
for $1/2\le\alpha<1$, the bounds
\[
 2\sin^2(\pi u)\le |1-\alpha e^{2\pi iu}|^2
                 =(1-\alpha)^2+4\alpha\sin^2(\pi u)\le4 .
\]
They give a periodic integrable majorant for the logarithm on each
period and hence a weighted integrable majorant on $[1,\infty)$.
We obtain
\begin{equation}\label{eq:factor}
 G(s)=-\frac{(s+1)(s+2)}{4\pi^2}H(s)\qquad(\R s>-2),
\end{equation}
where continuation to $\R s>-2$ follows by the identity theorem.

Let $\mathsf M,\mathsf N,\mathsf L\ge1$ be integers and $x_b=(b+\frac12)/\mathsf L$ for $0\le b\le\mathsf L-1$. Put
\begin{equation}\label{eq:KM}
\mathcal K_{\mathsf M}(x)=\sum_{m=1}^{\mathsf M}\frac{1-\cos(2\pi mx)}{m^3},\qquad H_{\mathsf M,\mathsf N}(s)=\sum_{n=1}^{\mathsf N}\int_0^1\mathcal K_{\mathsf M}(x)\,(n+x)^{-s-3}\,dx,
\end{equation}
\begin{equation}\label{eq:HMNL}
H_{\mathsf M,\mathsf N,\mathsf L}(s)=\frac1{\mathsf L}\sum_{n=1}^{\mathsf N}\sum_{b=0}^{\mathsf L-1}\mathcal K_{\mathsf M}(x_b)\,(n+x_b)^{-s-3}.
\end{equation}
Put also
\begin{equation}\label{eq:Bdefs}
B_0=\frac52,\qquad B_1=\frac{\pi^3}3,\qquad B_2=4\pi^2(1+\log\mathsf M),\qquad Z(\theta)=1+\frac1{\theta-1}\quad(\theta>1).
\end{equation}

For $\theta>1$ we have $\sum_{n\ge1}n^{-\theta}\le1+\int_1^\infty x^{-\theta}dx=Z(\theta)$. The functions $H_{\mathsf M,\mathsf N}$ and $H_{\mathsf M,\mathsf N,\mathsf L}$ are entire. By \eqref{eq:KM}, $0\le\mathcal K_{\mathsf M}\le B_0$, $|\mathcal K_{\mathsf M}'|\le2\pi\sum_{m\ge1}m^{-2}=B_1$ and $|\mathcal K_{\mathsf M}''|\le4\pi^2\sum_{m\le\mathsf M}1/m\le B_2$. We use the midpoint rule in the following form. If $\chi$ is a twice continuously differentiable complex function on $[0,1]$, then
\begin{equation}\label{eq:midpoint}
\Bigl|\int_0^1\chi(x)\,dx-\frac1{\mathsf L}\sum_{b=0}^{\mathsf L-1}\chi(x_b)\Bigr|\le\frac{\sup|\chi''|}{24\,\mathsf L^2}.
\end{equation}
Indeed, on the subinterval of length $1/\mathsf L$ centred at $x_b$, Taylor's formula gives $|\chi(x)-\chi(x_b)-\chi'(x_b)(x-x_b)|\le\frac12\sup|\chi''|(x-x_b)^2$, the linear term integrates to $0$, and $\int(x-x_b)^2dx=1/(12\mathsf L^3)$ over that subinterval.

\begin{theorem}[Explicit error bounds]\label{thm:trunc}
Let $\sigma=\operatorname{Re}s>-1$. Then $|H(s)-H_{\mathsf M,\mathsf N,\mathsf L}(s)|\le E_0$ and $|H'(s)-H_{\mathsf M,\mathsf N,\mathsf L}'(s)|\le E_1$, where
\begin{equation}\label{eq:E0expr}
\begin{aligned}
E_0={}&\frac1{(\sigma+2)\mathsf M^2}+\frac{B_0}{(\sigma+2)(\mathsf N+1)^{\sigma+2}}\\
&+\frac{B_2Z(\sigma+3)+2|s+3|B_1Z(\sigma+4)+|(s+3)(s+4)|B_0Z(\sigma+5)}{24\,\mathsf L^2}
\end{aligned}
\end{equation}
and
\begin{equation}\label{eq:E1expr}
\begin{aligned}
E_1={}&\frac1{(\sigma+2)^2\mathsf M^2}+\frac{B_0}{(\mathsf N+1)^{\sigma+2}}\Bigl(\frac{\log(\mathsf N+1)}{\sigma+2}+\frac1{(\sigma+2)^2}\Bigr)\\
&+\frac{B_2Z(\sigma+2)+2B_1\bigl(Z(\sigma+4)+|s+3|Z(\sigma+3)\bigr)}{24\,\mathsf L^2}\\
&+\frac{B_0\bigl(|(s+3)(s+4)|Z(\sigma+4)+|2s+7|Z(\sigma+5)\bigr)}{24\,\mathsf L^2}.
\end{aligned}
\end{equation}
\end{theorem}

\begin{proof}
By \eqref{eq:K} and \eqref{eq:KM}, $0\le\mathcal K-\mathcal K_{\mathsf M}\le2\sum_{m>\mathsf M}m^{-3}\le2\int_{\mathsf M}^\infty x^{-3}dx=\mathsf M^{-2}$. The difference $H-H_{\mathsf M,\mathsf N}$ is the integral of $(\mathcal K-\mathcal K_{\mathsf M})(u)u^{-s-3}$ over $1\le u\le\mathsf N+1$ plus the integral of $\mathcal K(u)u^{-s-3}$ over $u\ge\mathsf N+1$. With $\int_1^\infty u^{-\sigma-3}du=1/(\sigma+2)$, $\int_{\mathsf N+1}^\infty u^{-\sigma-3}du=(\mathsf N+1)^{-\sigma-2}/(\sigma+2)$ and $\mathcal K\le B_0$, this gives the first two terms of $E_0$. For the third, apply \eqref{eq:midpoint} to $\chi_n(x)=\mathcal K_{\mathsf M}(x)(n+x)^{-s-3}$ for $1\le n\le\mathsf N$. Its second derivative is $\mathcal K_{\mathsf M}''(n+x)^{-s-3}-2(s+3)\mathcal K_{\mathsf M}'(n+x)^{-s-4}+(s+3)(s+4)\mathcal K_{\mathsf M}(n+x)^{-s-5}$, and $|(n+x)^{-s-\theta}|\le n^{-\sigma-\theta}$ for $0\le x\le1$ and $\sigma+\theta>0$. Summing over $n$ with $\sum_nn^{-\theta}\le Z(\theta)$ gives the third term of $E_0$.

For the derivative, the integrand is $-\mathcal K(u)(\log u)u^{-s-3}$.
Integration gives, for $u_0\ge1$,
\[
 \int_{u_0}^\infty(\log u)u^{-\sigma-3}\,du
 =u_0^{-\sigma-2}
   \left(\frac{\log u_0}{\sigma+2}+\frac1{(\sigma+2)^2}\right).
\]
Taking $u_0=1$ and $u_0=\mathsf N+1$ gives the first two terms
of $E_1$. For the third, apply \eqref{eq:midpoint} to
$\chi_n(x)=-\mathcal K_{\mathsf M}(x)\log(n+x)(n+x)^{-s-3}$.
For $u\ge1$,
\begin{equation}\label{eq:logderivs}
\begin{aligned}
\frac{d}{du}\bigl[(\log u)u^{-s-3}\bigr]&=u^{-s-4}\bigl(1-(s+3)\log u\bigr),\\
\frac{d^2}{du^2}\bigl[(\log u)u^{-s-3}\bigr]&=u^{-s-5}\bigl((s+3)(s+4)\log u-(2s+7)\bigr).
\end{aligned}
\end{equation}
With $0\le\log u\le u$, the moduli of $(\log u)u^{-s-3}$ and of the two derivatives in \eqref{eq:logderivs} are at most $u^{-\sigma-2}$, $u^{-\sigma-4}+|s+3|u^{-\sigma-3}$ and $|(s+3)(s+4)|u^{-\sigma-4}+|2s+7|u^{-\sigma-5}$. These bounds decrease in $u$ because $\sigma>-1$, so on the $n$-th period they are at most their values at $u=n$. The product rule and $\sum_nn^{-\theta}\le Z(\theta)$ give the third term of $E_1$.
\end{proof}

\begin{lemma}[Second derivative]\label{lem:H2}
If $\operatorname{Re}s\ge\sigma>-2$ with $\sigma$ real, then $|H''(s)|\le2B_0/(\sigma+2)^3$.
\end{lemma}

\begin{proof}
By \eqref{eq:H} and $0\le\mathcal K\le B_0$, $|H''(s)|\le B_0\int_1^\infty(\log u)^2u^{-\sigma-3}du=2B_0/(\sigma+2)^3$.
\end{proof}

\begin{proposition}[Rouch\'e test for a disc]\label{prop:rouche}
Let $s_0\in\mathbb C$ and $\varsigma>0$ with $\operatorname{Re}s_0-\varsigma>-1$. Let $\eta_0=H_{\mathsf M,\mathsf N,\mathsf L}(s_0)$ and $\eta_1=H_{\mathsf M,\mathsf N,\mathsf L}'(s_0)$, and let $E_0,E_1$ be the bounds of Theorem~\ref{thm:trunc} at $s=s_0$. If
\begin{equation}\label{eq:rouche}
|\eta_0|+E_0+\varsigma E_1+\frac{\varsigma^2B_0}{(\operatorname{Re}s_0-\varsigma+2)^3}<\varsigma\,|\eta_1| ,
\end{equation}
then $H$ has exactly one zero in the disc $|s-s_0|<\varsigma$, and it is simple. It is also a simple zero of $G$, and $G$ has no other zero in the disc.
\end{proposition}

\begin{proof}
For $|s-s_0|=\varsigma$, Taylor's formula with integral remainder and Lemma~\ref{lem:H2} with $\sigma=\operatorname{Re}s_0-\varsigma$ give $|H(s)-H(s_0)-H'(s_0)(s-s_0)|\le\frac{\varsigma^2}2\cdot2B_0/(\operatorname{Re}s_0-\varsigma+2)^3$. By Theorem~\ref{thm:trunc}, $|H(s_0)-\eta_0|\le E_0$ and $|H'(s_0)-\eta_1|\le E_1$. Hence $|H(s)-\eta_1(s-s_0)|$ is at most the left side of \eqref{eq:rouche}, which is less than $\varsigma|\eta_1|=|\eta_1(s-s_0)|$. By Rouch\'e's theorem $H$ has in the disc as many zeros, counted with multiplicity, as the function $\eta_1(s-s_0)$, namely one. The disc lies in $\operatorname{Re}s>-1$, where the factor $(s+1)(s+2)$ of \eqref{eq:factor} does not vanish.
\end{proof}

The accompanying Rouch\'e certificate uses Arb \cite{johansson}
at 128 bits with $\mathsf M=4000$, $\mathsf N=200$, $\mathsf L=65536$
and radius $1/100$. Each disc below contains exactly one zero of $G$,
counted with multiplicity. The last column gives upper bounds
rounded upward from the certified ratios of the two sides of
\eqref{eq:rouche}.
\begin{center}
\small
\begin{tabular}{@{}rrrr@{}}
\toprule
Disc & $\R s_0$ & $\I s_0$ & Rouch\'e ratio, upper bound\\
\midrule
1 & 1.127312874 & 15.001263432 & 0.045205\\
2 & 0.533572290 & 21.763254707 & 0.145659\\
3 & 0.894730810 & 25.624836921 & 0.155103\\
4 & 0.227688021 & 30.975962723 & 0.466950\\
5 & 0.950506780 & 33.599087545 & 0.320245\\
6 & 0.430899874 & 38.054886210 & 0.452453\\
\bottomrule
\end{tabular}
\end{center}
The interval data and the program and log that produce them
accompany the paper.

\subsection{Nonvanishing of \texorpdfstring{$A$}{A}}
We use the digamma series \cite[(5.7.6)]{dlmf},
\begin{equation}\label{eq:psi}
 \psi(z)=-\gamma+\sum_{l\ge0}\left(\frac1{l+1}-\frac1{l+z}\right)
 \qquad(\R z>0).
\end{equation}
In \eqref{eq:Aseries} write $\Sigma_{2k}=1+(\Sigma_{2k}-1)$
and put
\begin{equation}\label{eq:Q0series}
 Q_0(s)=\sum_{k\ge1}\frac{\Sigma_{2k}-1}{k(2k-s)}\qquad(\R s<2).
\end{equation}
The identity $s/[k(2k-s)]=1/(k-s/2)-1/k$ and \eqref{eq:psi}
give
\begin{equation}\label{eq:sA}
 sA(s)=\log(2\pi)+\frac1s-\psi\Bigl(1-\frac s2\Bigr)-\gamma
       +sQ_0(s)\qquad(\R s<2,\ s\ne0).
\end{equation}

\begin{lemma}[A lower bound for the real part of the digamma function]\label{lem:repsi}
For real $x$ with $\frac12\le x\le1$ and real $y$,
\begin{equation}\label{eq:repsi}
\operatorname{Re}\psi(x+iy)\ge\frac12\log(x^2+y^2)-\gamma-\log2+\frac{y^2}{2x\,(x^2+y^2)} .
\end{equation}
\end{lemma}

\begin{proof}
By \eqref{eq:psi}, $\operatorname{Re}\psi(x+iy)-\psi(x)=\sum_{l\ge0}g_\psi(l)$ with $g_\psi(\lambda)=y^2/((\lambda+x)((\lambda+x)^2+y^2))$ for $\lambda\ge0$. With $v=\lambda+x>0$ one finds
\begin{equation}\label{eq:convex}
\frac{d^2}{dv^2}\,\frac1{v\,(v^2+y^2)}=\frac{12v^4+6v^2y^2+2y^4}{v^3(v^2+y^2)^3}>0,
\end{equation}
so $g_\psi$ is convex. For a convex function the chord lies above the graph, so $\int_l^{l+1}g_\psi\le\frac12(g_\psi(l)+g_\psi(l+1))$ for every integer $l\ge0$. Summing over $l$, and using $g_\psi(\lambda)\to0$ as $\lambda\to\infty$, gives $\int_0^\infty g_\psi\le\sum_{l\ge0}g_\psi(l)-\frac12g_\psi(0)$. Since $\int_0^\infty g_\psi=\int_x^\infty y^2\,dv/(v(v^2+y^2))=\frac12\log(1+y^2/x^2)$, we obtain $\operatorname{Re}\psi(x+iy)\ge\psi(x)+\frac12\log(1+y^2/x^2)+y^2/(2x(x^2+y^2))$. Termwise differentiation of \eqref{eq:psi} for real $x>0$ gives $\frac{d}{dx}(\psi(x)-\log x)=\sum_{l\ge0}(l+x)^{-2}-\frac1x>0$, because $(l+x)^{-2}>\int_l^{l+1}(\lambda+x)^{-2}d\lambda$ and the integrals add up to $\frac1x$. By \eqref{eq:psi}, $\psi(\frac12)=-\gamma-2\sum_{l\ge0}\bigl(\frac1{2l+1}-\frac1{2l+2}\bigr)=-\gamma-2\log2$. Hence for $x\ge\frac12$ we have $\psi(x)\ge\log x+\psi(\frac12)+\log2=\log x-\gamma-\log2$, and $\log x+\frac12\log(1+y^2/x^2)=\frac12\log(x^2+y^2)$ gives \eqref{eq:repsi}.
\end{proof}

\begin{theorem}[$A$ does not vanish above height $9$]\label{thm:Anonzero}
Let $0\le\sigma\le1$ and $|t|\ge9$, and put
\begin{equation}\label{eq:Qbound}
\mathcal Q(t)=\log(8\pi)-\frac12\log(1+t^2)-\frac{t^2}{2(t^2+4)}+\frac1{t^2}-0.64\,\frac{t^2-1}{t^2+4}+10^{-27} .
\end{equation}
Then $\operatorname{Re}\bigl(sA(s)\bigr)\le\mathcal Q(t)<0$ for $s=\sigma+it$. In particular $A(s)\ne0$.
\end{theorem}

\begin{proof}
Since $A(\bar s)=\overline{A(s)}$ by \eqref{eq:Aseries}, we may assume $t\ge9$. We bound the real parts of the terms of \eqref{eq:sA}. First, $\operatorname{Re}(1/s)=\sigma/(\sigma^2+t^2)\le1/t^2$. Second, apply Lemma~\ref{lem:repsi} with $x=1-\frac\sigma2\in[\frac12,1]$ and $y=-\frac t2$. Then $x^2+y^2\ge\frac14(1+t^2)$, and $x\le1$, $x^2+y^2\le\frac14(4+t^2)$ give $y^2/(2x(x^2+y^2))\ge t^2/(2(t^2+4))$. Hence
\begin{equation}\label{eq:psibound}
-\operatorname{Re}\psi\Bigl(1-\frac s2\Bigr)-\gamma\le2\log2-\frac12\log(1+t^2)-\frac{t^2}{2(t^2+4)} .
\end{equation}
Third, by \eqref{eq:Q0series}, $\operatorname{Re}(sQ_0(s))=\sum_{k\ge1}\frac{\Sigma_{2k}-1}k\,\mathfrak r_k$ with
\begin{equation}\label{eq:rhok}
\mathfrak r_k=\operatorname{Re}\frac{s}{2k-s}=\frac{2k\sigma-\sigma^2-t^2}{(2k-\sigma)^2+t^2} .
\end{equation}
For $0\le\sigma\le1$ we have $2k\sigma-\sigma^2\le2k-1$, because the left side increases with $\sigma$ on $[0,1]$. For $k=1$ the numerator is at most $1-t^2<0$ and the denominator is at most $4+t^2$, so $\mathfrak r_1\le-(t^2-1)/(t^2+4)$. If $2\le k$ and $2k-1<t^2$, then $\mathfrak r_k<0$. If $2k-1\ge t^2$, then $k\ge41$ because $t^2\ge81$, and $\mathfrak r_k\le(2k-1)/(2k-1)^2=1/(2k-1)$, since the denominator is at least $(2k-1)^2$. With $\Sigma_{2k}-1\le3\cdot4^{-k}$ these terms add up to at most $\frac{3}{41\cdot81}\sum_{k\ge41}4^{-k}=4^{-40}/3321<10^{-27}$. Finally $\Sigma_2-1>\sum_{m=2}^{10}m^{-2}+\int_{11}^\infty\lambda^{-2}d\lambda>0.5497+0.0909>0.64$. Adding up, $\operatorname{Re}(sA(s))\le\mathcal Q(t)$.

Each of the terms $-\frac12\log(1+t^2)$, $-\frac{t^2}{2(t^2+4)}=-\frac12+\frac2{t^2+4}$, $\frac1{t^2}$ and $-0.64\frac{t^2-1}{t^2+4}=-0.64+\frac{3.2}{t^2+4}$ decreases for $t>0$, so $\mathcal Q(t)\le\mathcal Q(9)$ for $t\ge9$. Since $e^{3.2242}>8\pi$ and $e^{4.4066}<82$, we have $\log(8\pi)<3.2242$ and $\frac12\log82>2.2033$. For these bounds one can use
$\pi<355/113$, the lower Taylor sum through degree $20$ for
$e^{3.2242}$, and the upper bound for $e^{4.4066}$ obtained by
adding $\frac54(4.4066)^{31}/31!$ to the sum through degree $30$.
The last bound follows by bounding consecutive term ratios in
the tail by $1/5$. With $\frac{81}{170}>0.4764$, $\frac1{81}<0.0124$ and $0.64\cdot\frac{80}{85}>0.6023$ this gives $\mathcal Q(9)<3.2242-2.2033-0.4764+0.0124-0.6023+10^{-27}<-0.045$.
\end{proof}

For the bounded rectangle, the program
\texttt{A\_nonvanishing\_check.py} uses the holomorphic function
$s^2A(s)$, which has value $1$ at $0$ by \eqref{eq:Aseries}.
For $c_k=\Sigma_{2k}-1$, it encloses $c_k$ by the rational
sum $\sum_{m=2}^{40}m^{-2k}$ with a nonnegative error at most
$40^{1-2k}/(2k-1)$ by the integral test. On $\R s\le1$ the
part of \eqref{eq:Q0series} with $k>80$ has modulus at most
\[
 \frac{\sum_{k>80}3\cdot4^{-k}}{81\cdot161}
 =\frac{4^{-80}}{81\cdot161},
\]
using $|2k-s|\ge 2k-1$. Multiplying \eqref{eq:sA} by $s$ gives
the expression enclosed by the program on each box. The initial
boxes tile $0\le\R s\le1$, $0\le\I s\le102$. Subdivision replaces
each undecided box by four boxes covering it. A box is accepted
when at least one of the real and imaginary enclosures excludes
zero. The recorded run accepted $188{,}892$ boxes at a maximum
subdivision depth of $5$, with no undecided boxes. Its smallest
recorded lower bound for
$\max\{|\R(s^2A(s))|,|\I(s^2A(s))|\}$ was enclosed by
$[1.00128\cdot10^{-5}\ \pm\ 3.76\cdot10^{-11}]$.
Thus $A$ is nonzero on the certified rectangle where $s\ne0$.
Together with the analytic estimate and conjugation symmetry,
this gives nonvanishing on $0\le\R s\le1$ wherever $A$ is
holomorphic.

\clearpage
\section{Figures}\label{app:figures}
The surface and zero-map panels are numerical
illustrations. The six discs in Appendix~\ref{app:cert} identify
which computed zeros of $G$ have enclosures.
The five modulus surfaces and the zero map were regenerated for
this version. Each red marker on a surface uses the computed modulus
at the corresponding mesh vertex.
On the real interval $s>-2$, \eqref{eq:H} gives $H(s)>0$,
since $\mathcal K$ is positive off the integers. Equation
\eqref{eq:factor} therefore gives the sign of $G$ shown in the
real-axis plot and a simple zero at $-1$.

For the period mean in the Gram plot, use the absolutely
convergent cosine series for $\log|1-\alpha e^{2\pi i u}|$
from the proof of \eqref{eq:factor}. The logarithmic majorant
used there is square integrable, so these functions converge
to $h$ in $L^2$ on a period. Their products at $x/j$ and $x/k$
converge in $L^1$ on a common period by Cauchy--Schwarz.
Only equal frequencies $m/j=n/k$ survive in the period mean,
with cosine product mean $1/2$. Writing
$(m,n)=r(j/g_{jk},k/g_{jk})$ with $r\ge1$ and $g_{jk}=\gcd(j,k)$, their
contribution is
\[
 \frac12\sum_{r\ge1}\frac{g_{jk}^2}{r^2jk}
   =\frac{\pi^2}{12}\frac{g_{jk}^2}{jk}.
\]
For the weighted panel the same $L^1$ convergence holds after
summing over periods against the weight $x^{-2\sigma-1}$ on
$[\max(j,k),\infty)$, since $\sigma>0$. Upon writing
$x=\max(j,k)v$, the cosine product formula replaces each
pair of modes by cosines with arguments
$2\pi\max(j,k)(m/j\pm n/k)v$. The kernel applied to each
frequency $\beta$ is $2\sigma\int_1^\infty
v^{-2\sigma-1}\cos(\beta v)\,dv$. It equals $1$ at zero and
has modulus at most $\min(1,4\sigma/|\beta|)$ otherwise, by
integration by parts. At fixed $m$, the difference-frequency
terms are bounded in sum by $O((1+\log m)/m^2)$: below
$mk/(2j)$ the denominator $|mk-nj|$ is $\gg m$, above
$2mk/j$ it is $\gg n$, and between them $n$ is comparable
to $m$ and the nonzero $mk-nj$ are spaced by $j$.
The equal-frequency terms contribute $O(m^{-2})$, and the
sum-frequency terms are bounded by the same estimate using
$mk+nj\gg m+n$. This proves absolute convergence, so the
frequency groups can be summed before truncating for the plot.

\clearpage
\sourcefigure{.95\linewidth}{fig_kernel.png}
 {The function $h$ on one period and the positive kernel
 $\mathcal K(u)=\sum_{m\ge1}(1-\cos(2\pi mu))/m^3$.}
\sourcefigure{.78\linewidth}{fig_G_real.png}{$G$ on the real axis.}

\clearpage
\sourcefigure{.87\linewidth}{fig_G_surface.png}
 {$|G(s)|$ on a linear scale with the vertical axis inverted.
 The red markers indicate the six computed zeros of $G$ in the certified discs.}
\sourcefigure{.87\linewidth}{fig_A_surface.png}
 {$|A(s)|$ on the same rectangle, with the vertical axis inverted.
 The nonvanishing result of Appendix~\ref{app:cert} covers the part
 with $0\le\R s\le1$.}

\clearpage
\sourcefigure{.96\linewidth}{fig_A_real_surface.png}
 {$|A(s)|$ near the real axis, clipped above $1.5$.
 The red markers indicate computed real zeros of $A$.}
\sourcefigure{.87\linewidth}{fig_GminusA_surface.png}
 {$|G(s)-A(s)|$, clipped above $0.6$, with the vertical axis inverted.}

\clearpage
\sourcefigure{.87\linewidth}{fig_zeta_surface.png}
 {$|\zeta(s)|$, clipped above $1.5$, with the vertical axis inverted.
 The red markers indicate the first six nontrivial zeros in the upper half-plane.}
\sourcefigure{.78\linewidth}{fig_zero_map.png}
 {Computed zeros of $G$, points with $G=A$, and arcs of
 $\I(G/A)=0$ with $0<\R(G/A)<1$ where $A\ne0$.
 The pairing shown by the arcs is a numerical observation.}

\clearpage
Let $z_b^G$ be the computed zeros of $G$ used in the plots and
$z_b^A$ the computed real zeros of $A$. Their counts are
$n_G=79$ and $n_A=101$. The exponential sums are
\begin{equation}\label{eq:expsum}
\mathcal E_G(x)=\frac1{n_G}\sum_{b=1}^{n_G}e^{ix\operatorname{Im}z^G_b},\qquad\mathcal E_A(x)=\frac1{n_A}\sum_{b=1}^{n_A}e^{ixz^A_b}.
\end{equation}
\sourcefigure{.70\linewidth}{fig_expsum_3d.png}
 {The exponential sums $\mathcal E_G(x)$ and $\mathcal E_A(x)$ for
 $0\le x\le2\pi$. The variable $x$ is the vertical axis.}
\sourcefigure{.98\linewidth}{fig_expsum_plane.png}
 {The two exponential sums in the complex plane and the computed zeros
 used in the sums, with their centroids.}

\clearpage
\sourcefigure{.95\linewidth}{fig_expsum_modulus.png}
 {Moduli of the exponential sums. The vertical lines in the upper panel
 are at $x=\log n$ for prime powers $n\le66$.}
\sourcefigure{.98\linewidth}{fig_gram.png}
 {Gram matrices of $h(x/k)$ for $1\le k\le10$.
 The first uses the mean over a common period. For the second, the
 integral with weight $x^{-2\sigma-1}$ starts at $\max(j,k)$ and
 is normalized by $2\sigma\max(j,k)^{2\sigma}$, with $\sigma=3/4$.}

\clearpage
\section{The vertical cut and the Python sources}\label{app:python}
The interactive vertical cut in \texttt{G\_A\_zeta\_cut\_v2.html}
shows $t\mapsto G(\tfrac12+it)$, $t\mapsto A(\tfrac12+it)$ and
their difference. On the sampled range $0.5\le t\le100$ with step
$0.01$, the accompanying output
reports 30 local minima, of which 29 met the numerical threshold
$|G-A|\le10^{-6}$ after refinement. The ordinates of these 29
points agree with the corresponding computed zeta ordinates within
$6\cdot10^{-7}$. The threshold and agreement are finite numerical
observations. The HTML, the script, and its output accompany this
paper.

The following Python sources generate the static figures and the data
for the vertical cut. The accompanying directory supplies
\texttt{cp\_eval.py} and the zero lists used by the plotting programs.
Run the four static plotting scripts in the order printed below.
The first generates the kernel and real-axis plots.
The vertical-cut listing uses a relative path to the evaluator.

\lstinputlisting[caption={\texttt{figures\_kernel\_real.py}}]{castor_v5/figures_kernel_real.py}
\lstinputlisting[caption={\texttt{figure\_gram.py}}]{castor_v5/figure_gram.py}
\lstinputlisting[caption={\texttt{figures\_v4.py}}]{castor_v5/figures_v4.py}
\lstinputlisting[caption={\texttt{expsum\_v3.py}}]{castor_v5/expsum_v3.py}
\lstinputlisting[caption={\texttt{vertical\_cut\_portable.py}}]{castor_v5/cut_G_A_zeta/vertical_cut_portable.py}

\end{document}
